% English abstract.

\begin{abstracten}

As long-context inference expands, the KV cache in autoregressive decoding grows with sequence length, shifting the deployment bottleneck from computation toward cache storage and memory access. Reducing cache bit-width lowers memory footprint, but tensor-level reconstruction error does not fully capture functional shifts after the quantized cache enters attention. Key perturbations can reshape logits and softmax attention distributions, while Value perturbations mainly affect output representations through weighted aggregation. This thesis proposes a behavior-aligned KV-cache quantization framework for efficient inference, using attention distributions, aggregation outputs, and task behavior as shared calibration and auditing targets across quantization-parameter selection, low-bit recovery, and layer-wise budget allocation.

Using calibration samples, this thesis evaluates full-precision reference and candidate quantized configurations in parallel. It records attention KL divergence and aggregation-output perturbation as fidelity readings for calibration-parameter selection, while retaining downstream task scores as audit signals. Layer-wise readings are then summarized into behavioral sensitivity profiles and assembled across models into a sensitivity-profile map, which exposes critical-layer distributions under different model conditions. This map supports fixed-$k$ comparisons between behavior-guided allocation and positional heuristics, and lets \texttt{AutoK} generate model-level $k$ candidates from sensitivity rankings and coverage thresholds. Thus, one set of offline behavioral readings supports fidelity calibration, low-bit behavior recovery, cross-model structural interpretation, and budget candidate generation.

Experiments cover six open-source instruction models from the Qwen2.5, LLaMA-3.1, and Mistral families, and show clear family-, scale-, and task-dependent behavior in KV-cache quantization and budget allocation. On the Qwen2.5-1.5B task-core protocol, the INT8 baseline differs from \texttt{FP16} by about $+0.02$ in the three-task mean, showing stable fidelity under conservative bit-width. With \texttt{INT4}, symmetric quantization causes stepwise retrieval collapse on Qwen models; K/V diagnostics identify Key-side low-bit noise as the direct trigger, and \texttt{INT4-RoleAlign} restores usable retrieval. Cross-model allocation maps four regimes: \texttt{AutoK} is clearest on Mistral-7B, Qwen2.5-3B favors early-layer protection, LLaMA-3.1-8B forms a medium-scale consensus regime, and Qwen2.5-14B shows nearby high-performance clusters without a stable winner. The \texttt{INT4} path yields about $73.4\%$ KV-cache capacity compression, while fused-decode time gains depend on $H_{kv}$, sequence length, and backend implementation. Overall, this thesis validates a behavior-aligned KV-cache quantization framework centered on attention-behavior preservation, K/V role-aware low-bit recovery, and \texttt{AutoK} with behavioral sensitivity profiles as the budget interface.

\keywordsen{large language models, KV-cache quantization,
attention-behavior preservation, low-bit inference,
behavior-guided calibration, budget allocation}

\end{abstracten}
